\begin{abstract}
This project designs and implements a real-time ray marching engine on the
Xilinx PYNQ-Z1 SoC FPGA, serving as the rendering backend for a Virtual
Reality headset with live head-tracking. Ray marching was selected for its
pixel-level independence, making it naturally suited to a parallel hardware
implementation.

The programmable logic implements a custom 27-bit floating-point library in
SystemVerilog, covering standard arithmetic, inverse square root, division,
and modulus, functioning analogously to a dedicated FPU. These primitives
drive a fully pipelined SDF evaluation accelerator feeding a dual-core ray
marching engine. Completed frames are streamed to DDR3 over AXI VDMA into a
double-buffered framebuffer, generating a real-time HDMI display at
$1024\times600$ at 18--40\,fps depending on scene complexity.

The ARM Cortex-A9 runs a Python interface that renders a scene-selection
overlay onto HDMI and coordinates runtime bitstream switching across nine
scenes. An MPU-6050 IMU communicates over I\textsuperscript{2}C to track
pitch, roll, and yaw, updating the camera lookat matrix every frame.

Audio reactivity is implemented via a per-frame FFT on a host PC. Bass,
midrange, treble energy, RMS, spectral centroid, and a beat detection flag
are streamed over TCP to the PYNQ, where they are written directly into
AXI-Lite SDF parameter registers; modulating spatial repetition period,
geometry scale, and camera orbit in real time with the music.

The accelerator exceeds a reference C implementation by over an order of
magnitude, satisfying the project brief. Nine SDF scenes are demonstrated
across analytic primitives and fractals. A neural implicit SDF inference
engine was designed and partially implemented; full integration remains
future work.

\bigskip
\vspace{2cm}
\textit{Throughout the report are scattered photos of the renders we were
able to obtain, or the renders we obtained with OpenGL GLSL shading language. As such, the reader travels with us in the beautiful world
of ray marching, seeing our initial attempts and failures as much as our
stunning final renders.}
\end{abstract}
